\section{Related work and novelty}\label{sec:related}

The problem sits at the meeting point of three lines of work: exactly divergence-free elements on curved domains, weak imposition of boundary conditions for Stokes flow, and the stability theory of Scott--Vogelius elements near singular and nearly singular vertices. We discuss each and then state what is new.

\subsection{The prize problem and the constant-pressure case}
The question we answer was posed by Scott \cite{ScottPrize}, on the basis of the computations of Gjerde and Scott \cite{GjerdeScott2024}. Cavalcante \cite{Cavalcante} has analysed the constant-pressure (zero-pressure-gradient) benchmark of the prize: the abstract of that preprint states the bound $\norm{u-u_h}_{H^1}\le C(\hO^k+\hG^{3/2})$ for this benchmark with quartic velocities on an inscribed polygon, and it reports a coercivity transition of the penalty near $\mu\approx60.338$ for its test mesh; our code gives $\mu^\ast=60.343$ at the same resolution ($N=20$ chords). In our notation that is the case $G=0$ of the printed method, which we recover in Theorem~\ref{thm:A}; our proofs do not use \cite{Cavalcante}. The present paper treats general Stokes data. The difference is essential: once the wall pressure varies, the printed method no longer has Scott's rate (Theorems~\ref{thm:A}--\ref{thm:B}), and a corrected method is needed. The analysis of the constant-pressure case and ours are thus complementary answers to the two halves of Scott's question: the ``if'' (the estimate holds when the pressure is constant on the wall) and the ``why not'' (it fails otherwise, for an identifiable reason).

\subsection{Divergence-free elements on curved domains}
Neilan and Otus \cite{NeilanOtus2021} constructed Scott--Vogelius elements on curved two-dimensional domains, with boundary elements mapped to fit the geometry, and proved convergence estimates; they also showed how to recover pressure robustness by reconstructing the load with commuting interpolants \cite[Theorem~6.1]{NeilanOtus2021}. Durst and Neilan \cite{DurstNeilan2024} extended the approach to general degree on smooth planar domains. In three dimensions, Chen and Liu \cite{ChenLiu2026} obtain pressure-robust, optimal-order $h^k$ estimates with isoparametric Piola-mapped elements on curved domains. Methods that fit the boundary can reach rates above $\hG^{3/2}$ (Chen and Liu attain $h^k$); the inscribed-polygon Nitsche route analysed here is the low-cost, straight-sided alternative, whose best possible rate is $\hG^{3/2}+h^k$ (Theorem~\ref{thm:lb}). We therefore do not regard it as a competitor on rate, but as the analysis of a method that is easy to implement in any code that already has Scott--Vogelius elements.

The work closest to ours is that of Liu, Neilan and Otus \cite{LiuNeilanOtus2023}. They use Scott--Vogelius elements on Clough--Tocher splits of a mesh lying inside a curved domain, impose the tangential condition by a Nitsche-type form with a Taylor boundary correction, and enforce the normal component by a separate mean-free boundary multiplier which, as they note, approximates the wall pressure up to a constant. Together with a zero-flux constraint on the velocity space the method is exactly divergence-free and converges at the optimal order $h^k$; without the constraint, they remark, it is ill posed \cite[Remark~3.2]{LiuNeilanOtus2023}. Our Proposition~\ref{prop:singular} is the analogue of this remark for the square Scott--Vogelius system. Our subject is different: the method as Gjerde and Scott printed it, and its minimal repair, with no extra unknowns, no boundary correction and the pressure itself on the wall.

For cut (unfitted) meshes, Liu, Neilan and Olshanskii \cite{LiuNeilanOlshanskii2023} place the boundary pressure term in both the momentum and the continuity equation; the discrete divergence then vanishes only away from an $O(h)$ strip at the boundary. Frachon, Nilsson and Zahedi \cite{FrachonNilssonZahedi2024} place it in the momentum equation only, so as not to perturb the divergence condition, and show that fixing the pressure mean by a Lagrange multiplier makes the discrete divergence equal to that multiplier \cite[(5.7)--(5.9)]{FrachonNilssonZahedi2024}. They avoid this by testing the momentum equation only with velocities of zero net boundary flux. Our $\CNS$ is this device in square form (Remark~\ref{rem:fnz}); we claim no novelty for it. Burman, Hansbo and Larson give related multiplier formulations for divergence-free cut elements \cite{BurmanHansboLarson2024}. Exactly divergence-free methods also require discrete Dirichlet data with zero net flux, as Eickmann, Scott and Tscherpel point out for polygonal domains \cite{EickmannScottTscherpel2025}; our flux-corrected outer data (Definition~\ref{def:flux}) are a simple instance.

\subsection{Nitsche's method for Stokes flow}
Weak imposition of Dirichlet data in Nitsche's sense was developed further by Freund and Stenberg \cite{FreundStenberg1995} and Stenberg \cite{Stenberg1995}; for Stokes flow the consistent form uses the full traction and contains the pressure term \cite{FreundStenberg1995,Stenberg1995,Becker2002,BurmanHansbo2014}. Penalty-free variants \cite{Burman2012,BoiveauBurman2016} and Robin-type generalisations \cite{JuntunenStenberg2009} exist. That the pressure term is needed for consistency is known \cite{BurmanHansboLarson2019Stokes}, and that the Nitsche penalty must scale with the local boundary mesh size follows from the trace-inverse inequality \cite{Stenberg1995,JuntunenStenberg2009}. What has not been done, to our knowledge, is an analysis of a Nitsche--Stokes form \emph{without} the pressure term: whether it converges, at what rate, and what its error looks like. The mechanism, a consistency error of penalty type acting only on the pressure, is that of Babu\v{s}ka's penalty method \cite{Babuska1973}.

On the geometric side, the polygonal approximation of curved boundaries for scalar problems is classical \cite{BergerScottStrang1972,Scott1975}, and higher-order boundary corrections on polygons are due to Bramble, Dupont and Thom\'ee \cite{BrambleDupontThomee1972} and, more recently, Dupont, Guzm\'an and Scott \cite{DupontGuzmanScott2025} (for the Poisson problem). Bernardi, Costabel, Dauge and Girault \cite{BCDG2016} show that the continuous inf-sup constant of an inscribed polygonal approximation converges to that of the curved domain; we use a simple gluing argument for the same purpose. None of the cut-element or boundary-correction works above uses exactly divergence-free spaces on an inscribed polygon.

\subsection{Singular and nearly singular vertices}
The stability of Scott--Vogelius elements degenerates near singular vertices, and the Guzm\'an--Scott quantity $\Theta(z)$ \cite{GuzmanScott2019} measures this. Gr\"a\ss le, Bohne and Sauter \cite{GrassleBohneSauter2024} introduced pressure constraints (the ``pressure-wired'' element) that restore a shape-robust inf-sup constant at nearly singular vertices, and spurious pressures with robust velocities, together with postprocessing filters, were studied by Park \cite{Park2020} and by Bohne, Gr\"a\ss le and Sauter \cite{BohneGrassleSauter2025}. Our constrained inf-sup result for strong imposition (Theorem~\ref{thm:st:locks}(a)) is essentially the wall-vertex case of theirs, and our pressure filter is in the same spirit. Kean, Neilan and Schneier \cite{KeanNeilanSchneier2022} analyse Scott--Vogelius elements on anisotropic meshes obtained by barycentric and incenter refinement. What is different in our setting is that the loss at a locked wall vertex is an \emph{approximation} loss, caused by the zero trace on the polygon, and not a stability loss; the velocity error can therefore be characterised two-sidedly.

The zero-gradient lock at a wall vertex with two triangles, and the statement that no constraint arises with more than two triangles, are due to Scott \cite[Theorem~1]{ScottPrize}; Gjerde and Scott treat the two-triangle case in \cite[\S3.4, App.~B]{GjerdeScott2024}.

\subsection{Broader context}
Several further strands of the literature enter the analysis or its extensions. The divergence-free pairs used in three dimensions, relevant to the companion paper \cite{PadhiThreeD}, are stable on Alfeld (barycentric) splits \cite{Zhang2005,GuzmanNeilan2018}, on Worsey--Farin splits \cite{GuzmanLischkeNeilan2022}, and on Freudenthal meshes for large degree \cite{Zhang2011}; smooth discrete de Rham complexes \cite{Neilan2015,FuGuzmanNeilan2020} provide the corresponding smooth stream-function spaces, and Farrell, Mitchell and Scott \cite{FarrellMitchellScott2024} survey what is known and conjectured. A local Fortin projection for Scott--Vogelius elements on general meshes has recently been constructed by Eickmann, Guzm\'an, Neilan, Scott and Tscherpel \cite{EickmannGNST2026}. Our $L^2$ estimates on graded meshes use local energy estimates in the spirit of Arnold and Liu \cite{ArnoldLiu1995} and of Demlow, Guzm\'an and Schatz \cite{DemlowGuzmanSchatz2011}, and the regularity result of Kellogg and Osborn \cite{KelloggOsborn1976} on convex polygons. The Navier--Stokes extension follows the framework of nonsingular solution branches of Brezzi, Rappaz and Raviart \cite{BRR1980,GiraultRaviart}, with a duality argument of Schatz type \cite{Schatz1974} for explicit thresholds. Sharp trace-inverse constants for polynomials \cite{WarburtonHesthaven2003} enter the analysis of the pressure in the first row of boundary triangles. For nearly singular wall stars we use anisotropic interpolation \cite{BabuskaAziz1976,Apel1999} and right inverses of the divergence on John domains \cite{AcostaDuranMuschietti2006}.

\subsection{What is known and what is new}
\emph{Known, and used or recovered here.}
\begin{itemize}
\item The pressure term is necessary for consistency of Nitsche's method for Stokes flow \cite{BurmanHansboLarson2019Stokes}; penalty consistency errors of the leak type are classical \cite{Babuska1973}.
\item The mean-free (zero-net-flux) placement of the pressure term that keeps the discrete velocity divergence-free \cite{FrachonNilssonZahedi2024}; the need for a flux constraint or correction \cite{LiuNeilanOtus2023,EickmannScottTscherpel2025}.
\item Sufficiency of a penalty proportional to the ratio of bulk to boundary mesh size (the trace-inverse inequality \cite{Stenberg1995,JuntunenStenberg2009}).
\item The two-triangle lock \cite{ScottPrize,GjerdeScott2024}; constrained and filtered pressures at nearly singular vertices \cite{GrassleBohneSauter2024,Park2020,BohneGrassleSauter2025}.
\item The $H^1$ bound $\hO^k+\hG^{3/2}$ for the constant-pressure benchmark with quartic velocities \cite{Cavalcante}.
\end{itemize}
\emph{New in this paper}, as far as we know.
\begin{itemize}
\item A two-sided analysis of the printed Gjerde--Scott method for general data: convergence at the sharp rate $h/\mu$ in $H^1$ and $(h/\mu)^{1/2}$ in energy, the explicit leading error term with constant $1$ in the slip and energy norms, and the identification of the $H^1$ constant with the energy of a Stokes leak flow (Section~\ref{sec:printed}).
\item The proof of Scott's conjectured rate $\hG^{3/2}+\hO^k$ for general Stokes data, for the mean-free correction on an inscribed polygon, with a penalty-uniform $H^1$ bound under flux-corrected data; and a matching lower bound showing that the term $\hG^{3/2}$ is sharp, for every $k\ge4$ (Section~\ref{sec:corrected}).
\item The observation that a growing penalty restores the $H^1$ rate $\frac32$ of the printed method but can never lift its energy rate above $1$.
\item Explicit necessity statements for the penalty threshold, including exact rational certificates on shape families of boundary triangles, a proof on every shape-regular mesh, and exact counterexamples to purely local necessity (Section~\ref{sec:penalty}).
\item A two-sided theory of strongly imposed no-slip on the inscribed polygon: the counting lower bound at locked vertices, the sharp two-sided rate with locks, and the needle-angle law (Section~\ref{sec:strong}). The upper bound $\hG^{3/2}+h^k$ under the Guzm\'an--Scott condition is close to a corollary of \cite{GuzmanScott2019} and the scalar theory \cite{BergerScottStrang1972,Scott1975}; what is new there is only the observation that it holds, and that the lock therefore needs two-triangle or nearly singular wall stars.
\item The results surveyed in Section~\ref{sec:further} (pressure and drag, $L^2$, penalty-uniform leak limit, Navier--Stokes, general domains, nearly singular stars, pressure robustness), proved in \cite{PadhiExt}, some of them only under the additional hypotheses listed in Table~\ref{tab:status}.
\end{itemize}
